\allowdisplaybreaks
\begin{center}
	\hrule
	\vspace{.4cm}
 \begin{tabular*}{\textwidth}{@{}l@{}|@{\extracolsep{0.6in}}r@{}}%
\parbox{4.25in}{\raggedright{\includegraphics[width=.9\linewidth]{ictp-pwf.pdf}}} &
\parbox[c][]{4in}{{\Large\textbf{Md Akiful Islam Zawad} \par}
                    { Brac University \par}
                    { KHA 224, Progati Sarani, Merul Badda, \par}
                    { Dhaka 1212, Bangladesh \par}
                    { \href{ext.akiful.islam@bracu.ac.bd}{ext.akiful.islam@bracu.ac.bd}} \par}
\end{tabular*}\vspace{.3in}
\definecolor{ceruleanblue}{rgb}{0.16, 0.32, 0.6}
	\LARGE\scshape\textbf{\textcolor{ceruleanblue}{Physics For Bangladesh: School On Quantum Field Theory}}
\end{center}
\hrule\vspace{.25in}
{\large\textbf{Problem Sheet:}\ \textsc{1} \hspace{\hfill} \large\textbf{Submission Date:} \today\\
	\hline\hline
 %----------------------------
\paragraph*{Problem Set 3} %\hfill \newline
\\
Verify Wick's theorem for
    \begin{enumerate}
        \item The product of three scalars:
        \begin{align*}
            T(\phi(x_1) \phi(x_2) \phi(x_3)) = : \phi(x_1) \phi(x_2) \phi(x_3) : + \phi(x_1) \Delta_F(x_2 - x_3)
        + \phi(x_2) \Delta_F(x_3 - x_1) + \phi(x_3) \Delta_F(x_1 - x_2)
        \end{align*}
        \item The product of four scalars:
        \begin{align*}
            T(\phi(x_1) \phi(x_2) \phi(x_3) \phi(x_4))    
        \end{align*}
    \end{enumerate}

\bigskip\bigskip\hline\hline\bigskip
\section*{Solution}
Wick's theorem states that the time-ordered product of a set of fields can be expressed as a sum of normal-ordered products with contractions. For scalar fields, the contraction between two fields is defined as the Feynman propagator $\Delta_F(x - y)$.

Wick’s Theorem states that for $n$ field operators, the time-ordered product can be written as a sum over all possible ways of contracting the operators. 

For two scalar fields $\phi(x_1)\equiv\phi_1$ and $\phi(x_2)\equiv\phi_2$, the theorem is:

\begin{align}
    T \{\phi_1\phi_2\} &= :\phi_1\phi_2: + \overbrace{\phi_1,\phi_2}
\end{align}
The time-ordered product of two scalar field is defined as
\begin{align*}
    T \{\phi_1\phi_2\} &= \begin{cases}
      \phi_1\phi_2; & x_1^0 > x_2^0\\
      \phi_2\phi_1; & x_2^0 > x_1^0
    \end{cases}
\end{align*}
or \begin{align*}
    T \{\phi_1\phi_2\}=\theta\left(x_1^0-x_2^0\right)\phi_1\phi_2+\theta\left(x_2^0-x_1^0\right)\phi_2\phi_1,
\end{align*}
from which it is clear that $T \{\phi_1\phi_2\}$ is symmetric in $x_1$ and $x_2$.

The normal-ordered product of two fields is defined as
\begin{align*}
    :\phi_1\phi_2:=:\left(\phi_1^++\phi_1^-\right)+\left(\phi_2^-+\phi_2^-\right):=\phi_1^+\phi_2^++\phi_2^-\phi_1^++\phi_1^-\phi_2^++\phi_1^-\phi_2^-
\end{align*}
The Feynmann propagator may be written as
\begin{align}
    \Delta_F\left(x_1-x_2\right)=T\{\phi_1\phi_2\}-:\phi_1\phi_2:
\end{align}
We now attempt to check Wick's theorem fir the time-ordered product of three scalar fields. First, consider the case, $x_1^0>x_2^0>x_3^0$:
\begin{align}
    T\left\{\phi_1\phi_2\phi_3\right\}&=\phi_1\phi_2\phi_3=\left(\phi_1^++\phi_1^-\right)\left(\phi_2^++\phi_2^-\right)\left(\phi_3^++\phi_3^-\right)\notag\\
    &=\phi_1^+\phi_2^+\phi_3^++\phi_1^+\phi_2^+\phi_3^-+\phi_1^+\phi_2^-\phi_3^++\phi_1^+\phi_2^-\phi_3^-+\phi_1^-\phi_2^+\phi_3^++\phi_1^-\phi_2^+\phi_3^-+\phi_1^-\phi_2^-\phi_3^++\phi_1^-\phi_2^-\phi_3^-\notag\\
    &=\phi_1^+\phi_2^+\phi_3^+\left(\phi_3^-\phi_1^+\phi_2^++\left[\phi_1^+,\phi_3^-\right]\phi_2^++\phi_1^+\left[\phi_2^+,\phi_3^-\right]\right)+\left(\phi_2^-\phi_2^+\phi_3^++\left[\phi_1^+,\phi_2^-\right]\phi_3^+\right)\notag\\
    &\hspace{20pt}\left(\phi_2^-\phi_3^-\phi_1^++\phi_2^-\left[\phi_1^+,\phi_3^-\right]+\left[\phi_1^+,\phi_2^-\right]\phi_3^-\right)+\phi_1^-\phi_2^+\phi_3^+\notag\\
    &\hspace{40pt}\left(\phi_1^-\phi_2^+\phi_3^-+\phi_1^-\left[\phi_1^+,\phi_3^-\right]\right)+\phi_1^-\phi_3^+\phi_2^-+\phi_1^-\phi_2^-\phi_3^-\notag\\
    &=:\phi_1\phi_2\phi_3:+\left[\phi_1^+,\phi_3^-\right]\phi_2+\left[\phi_2^+,\phi_3^-\right]\phi_1+\left[\phi_1^+,\phi_3^-\right]\phi_3\notag\\
    &=:\phi_1\phi_2\phi_3:+\Delta_F\left(x_1-x_3\right)\phi_2+\Delta_F\left(x_2-x_3\right)\phi_1+\Delta_F\left(x_1-x_2\right)\phi_3\notag
\end{align}
As a re-ordering of $x_1^0>x_2^0>x_3^0$ corresponds to a different choice of labels and since $\Delta_F$ and normal ordered products are symmetric, this calculation suffices to prove Wick’s theorem.

Thus, Wick's theorem is verified for the product of three scalar fields.

Now, we verify Wick's theorem for the product of four scalar fields:
\begin{align}
    T\left\{\phi_1 \phi_2 \phi_3 \phi_4\right\} &= \left(\phi_1^+ + \phi_1^-\right)\left(\phi_2^+ + \phi_2^-\right)\left(\phi_3^+ + \phi_3^-\right)\left(\phi_4^+ + \phi_4^-\right) \notag \\
    &= \phi_1^+ \phi_2^+ \phi_3^+ \phi_4^+ + \phi_1^+ \phi_2^+ \phi_3^+ \phi_4^- + \phi_1^+ \phi_2^+ \phi_3^- \phi_4^+ + \phi_1^+ \phi_2^+ \phi_3^- \phi_4^- \notag \\
    &\quad + \phi_1^+ \phi_2^- \phi_3^+ \phi_4^+ + \phi_1^+ \phi_2^- \phi_3^+ \phi_4^- + \phi_1^+ \phi_2^- \phi_3^- \phi_4^+ + \phi_1^+ \phi_2^- \phi_3^- \phi_4^- \notag \\
    &\quad + \phi_1^- \phi_2^+ \phi_3^+ \phi_4^+ + \phi_1^- \phi_2^+ \phi_3^+ \phi_4^- + \phi_1^- \phi_2^+ \phi_3^- \phi_4^+ + \phi_1^- \phi_2^+ \phi_3^- \phi_4^- \notag \\
    &\quad + \phi_1^- \phi_2^- \phi_3^+ \phi_4^+ + \phi_1^- \phi_2^- \phi_3^+ \phi_4^- + \phi_1^- \phi_2^- \phi_3^- \phi_4^+ + \phi_1^- \phi_2^- \phi_3^- \phi_4^- \notag \\
    &= :\phi_1 \phi_2 \phi_3 \phi_4: \notag \\
    &\quad + \left(\phi_1^+ \phi_2^+ \left[\phi_3^-, \phi_4^+\right] + \phi_1^+ \left[\phi_2^+, \phi_4^-\right] \phi_3^+ + \left[\phi_1^+, \phi_3^-\right] \phi_2^+ \phi_4^+ \right) \notag \\
    &\quad + \left(\phi_2^+ \phi_3^+ \left[\phi_1^-, \phi_4^+\right] + \phi_2^+ \left[\phi_1^+, \phi_4^-\right] \phi_3^+ + \left[\phi_2^+, \phi_3^-\right] \phi_1^+ \phi_4^+ \right) \notag \\
    &\quad + \left(\phi_3^+ \phi_4^+ \left[\phi_1^-, \phi_2^+\right] + \phi_3^+ \left[\phi_1^+, \phi_2^-\right] \phi_4^+ + \left[\phi_1^+, \phi_4^-\right] \phi_2^+ \phi_3^+ \right) \notag \\
    &\quad + \left(\phi_1^- \phi_2^- \left[\phi_3^+, \phi_4^-\right] + \phi_1^- \left[\phi_2^-, \phi_4^+\right] \phi_3^- + \left[\phi_1^-, \phi_3^+\right] \phi_2^- \phi_4^- \right) \notag \\
    &\quad + \left(\phi_2^- \phi_3^- \left[\phi_1^+, \phi_4^-\right] + \phi_2^- \left[\phi_1^-, \phi_4^+\right] \phi_3^- + \left[\phi_2^-, \phi_3^+\right] \phi_1^- \phi_4^- \right) \notag \\
    &\quad + \left(\phi_3^- \phi_4^- \left[\phi_1^+, \phi_2^-\right] + \phi_3^- \left[\phi_1^-, \phi_2^+\right] \phi_4^- + \left[\phi_1^-, \phi_4^+\right] \phi_2^- \phi_3^- \right)
\end{align}
We now attempt to check Wick's theorem for the time-ordered product of four scalar fields. First, consider the case, $x_1^0>x_2^0>x_3^0x^4>0$:

Start with the product:
\begin{align}
    \phi_1 \phi_2 \phi_3 \phi_4 &= \left(\phi_1^+ + \phi_1^-\right)\left(\phi_2^+ + \phi_2^-\right)\left(\phi_3^+ + \phi_3^-\right)\left(\phi_4^+ + \phi_4^-\right) \notag \\
    &= \phi_1^+ \phi_2^+ \phi_3^+ \phi_4^++ \phi_1^+ \phi_2^+ \phi_3^+ \phi_4^-+ \phi_1^+ \phi_2^+ \phi_3^- \phi_4^++ \phi_1^+ \phi_2^+ \phi_3^- \phi_4^- \notag \\
    &\quad + \phi_1^+ \phi_2^- \phi_3^+ \phi_4^++\phi_1^+ \phi_2^- \phi_3^+ \phi_4^-+ \phi_1^+ \phi_2^- \phi_3^- \phi_4^++ \phi_1^+ \phi_2^- \phi_3^- \phi_4^-+ \notag\\ &\quad +\phi_1^- \phi_2^+ \phi_3^+ \phi_4^++ \phi_1^- \phi_2^+ \phi_3^+ \phi_4^-+ \phi_1^- \phi_2^+ \phi_3^- \phi_4^++ \phi_1^- \phi_2^+ \phi_3^- \phi_4^- \notag \\
    &\quad + \phi_1^- \phi_2^- \phi_3^+ \phi_4^++ \phi_1^- \phi_2^- \phi_3^+ \phi_4^-+ \phi_1^- \phi_2^- \phi_3^- \phi_4^+ + \phi_1^- \phi_2^- \phi_3^- \phi_4^- \notag
\end{align}

% Now, apply the time-ordering operator \( T \) and use commutation relations to correct each term:

% \begin{align}
%     T\left\{\phi_1 \phi_2 \phi_3 \phi_4\right\} &= \phi_1^+ \phi_2^+ \phi_3^+ \phi_4^+ \notag \\
%     &\quad + T\left\{\phi_1^+ \phi_2^+ \phi_3^+ \phi_4^-\right\} \notag \\
%     &\quad + T\left\{\phi_1^+ \phi_2^+ \phi_3^- \phi_4^+\right\} \notag \\
%     &\quad + T\left\{\phi_1^+ \phi_2^+ \phi_3^- \phi_4^-\right\} \notag \\
%     &\quad + T\left\{\phi_1^+ \phi_2^- \phi_3^+ \phi_4^+\right\} \notag \\
%     &\quad + T\left\{\phi_1^+ \phi_2^- \phi_3^+ \phi_4^-\right\} \notag \\
%     &\quad + T\left\{\phi_1^+ \phi_2^- \phi_3^- \phi_4^+\right\} \notag \\
%     &\quad + T\left\{\phi_1^+ \phi_2^- \phi_3^- \phi_4^-\right\} \notag \\
%     &\quad + T\left\{\phi_1^- \phi_2^+ \phi_3^+ \phi_4^+\right\} \notag \\
%     &\quad + T\left\{\phi_1^- \phi_2^+ \phi_3^+ \phi_4^-\right\} \notag \\
%     &\quad + T\left\{\phi_1^- \phi_2^+ \phi_3^- \phi_4^+ \right\} \notag \\
%     &\quad + T\left\{\phi_1^- \phi_2^+ \phi_3^- \phi_4^- \right\} \notag \\
%     &\quad + T\left\{\phi_1^- \phi_2^- \phi_3^+ \phi_4^+ \right\} \notag \\
%     &\quad + T\left\{\phi_1^- \phi_2^- \phi_3^+ \phi_4^- \right\} \notag \\
%     &\quad + T\left\{\phi_1^- \phi_2^- \phi_3^- \phi_4^+ \right\} \notag \\
%     &\quad + T\left\{\phi_1^- \phi_2^- \phi_3^- \phi_4^- \right\} \notag
% \end{align}
Now expand and correct each time-ordered product term by term:
    \begin{align*}
        T\left\{\phi_1 \phi_2 \phi_3 \phi_4\right\} &= \phi_1^+ \phi_2^+ \phi_3^+ \phi_4^+\\[5pt]
        &\hspace{5pt}+\phi_1^+ \phi_2^+ \phi_3^+ \phi_4^-+ \left[\phi_4^-, \phi_1^+\right] \phi_2^+ \phi_3^++ \left[\phi_4^-, \phi_2^+\right] \phi_1^+ \phi_3^++ \left[\phi_4^-, \phi_3^+\right] \phi_1^+ \phi_2^+\\[5pt]
        &\hspace{10pt}+\phi_1^+ \phi_2^+ \phi_3^- \phi_4^++ \left[\phi_3^-, \phi_4^+\right] \phi_1^+ \phi_2^++ \left[\phi_1^+, \phi_3^- \right] \phi_2^+ \phi_4^++ \left[\phi_2^+, \phi_3^- \right] \phi_1^+ \phi_4^+\\[5pt]
        &\hspace{15pt}+\phi_1^+ \phi_2^+ \phi_3^- \phi_4^-+ \left[\phi_4^-, \phi_1^+\right] \phi_2^+ \phi_3^-+ \left[\phi_3^-, \phi_1^+\right] \phi_2^+ \phi_4^-+ \left[\phi_3^-, \phi_2^+\right] \phi_1^+ \phi_4^-\\[5pt]
        &\hspace{20pt}+\phi_1^+ \phi_2^- \phi_3^+ \phi_4^++ \left[\phi_2^-, \phi_3^+\right] \phi_1^+ \phi_4^++ \left[\phi_1^+, \phi_2^-\right] \phi_3^+ \phi_4^++ \left[\phi_1^+, \phi_4^+\right] \phi_2^- \phi_3^+\\[5pt]
        &\hspace{25pt}+\phi_1^+ \phi_2^- \phi_3^+ \phi_4^-+ \left[\phi_4^-, \phi_2^-\right] \phi_1^+ \phi_3^++ \left[\phi_3^+, \phi_2^-\right] \phi_1^+ \phi_4^-+ \left[\phi_1^+, \phi_2^-\right] \phi_3^+ \phi_4^-\\[5pt]
        &\hspace{30pt}+\phi_1^+ \phi_2^- \phi_3^- \phi_4^++\left[\phi_3^-, \phi_4^+\right] \phi_1^+ \phi_2^-+\left[\phi_2^-, \phi_4^+\right] \phi_1^+ \phi_3^-+\left[\phi_2^-, \phi_3^- \right] \phi_1^+ \phi_4^+\\[5pt]
        &\hspace{35pt}+\phi_1^+ \phi_2^- \phi_3^- \phi_4^-+\left[\phi_3^-, \phi_4^-\right] \phi_1^+ \phi_2^-+\left[\phi_2^-, \phi_4^-\right] \phi_1^+ \phi_3^-+\left[\phi_2^-, \phi_3^- \right] \phi_1^+ \phi_4^-\\[5pt]
        &\hspace{40pt}+\phi_1^- \phi_2^+ \phi_3^+ \phi_4^+ +\left[\phi_1^-, \phi_4^+\right] \phi_2^+ \phi_3^++\left[\phi_1^-, \phi_3^+\right] \phi_2^+ \phi_4^++\left[\phi_1^-, \phi_2^+\right] \phi_3^+ \phi_4^+\\[5pt]
        &\hspace{45pt}+\phi_1^- \phi_2^+ \phi_3^+ \phi_4^-+\left[\phi_4^-, \phi_1^-\right] \phi_2^+ \phi_3^++\left[\phi_3^+, \phi_1^-\right] \phi_2^+ \phi_4^-+\left[\phi_2^+, \phi_1^-\right] \phi_3^+ \phi_4^-\\[5pt]
        &\hspace{50pt}+\phi_1^- \phi_2^+ \phi_3^- \phi_4^++ \left[\phi_3^-, \phi_4^+\right] \phi_1^- \phi_2^++ \left[\phi_2^+, \phi_3^-\right] \phi_1^- \phi_4^++ \left[\phi_1^-, \phi_3^-\right] \phi_2^+ \phi_4^+\\[5pt]
        &\hspace{55pt}+\phi_1^- \phi_2^+ \phi_3^- \phi_4^-+ \left[\phi_3^-, \phi_4^-\right] \phi_1^- \phi_2^++ \left[\phi_2^+, \phi_4^-\right] \phi_1^- \phi_3^-+ \left[\phi_2^+, \phi_3^- \right] \phi_1^- \phi_4^-\\[5pt]
        &\hspace{60pt}+\phi_1^- \phi_2^- \phi_3^+ \phi_4^++ \left[\phi_3^+, \phi_4^+\right] \phi_1^- \phi_2^-+ \left[\phi_2^-, \phi_3^+\right] \phi_1^- \phi_4^++ \left[\phi_1^-, \phi_3^+\right] \phi_2^- \phi_4^+\\[5pt]
        &\hspace{65pt}+\phi_1^- \phi_2^- \phi_3^+ \phi_4^-+ \left[\phi_4^-, \phi_2^-\right] \phi_1^- \phi_3^++ \left[\phi_3^+, \phi_2^-\right] \phi_1^- \phi_4^-+ \left[\phi_1^-, \phi_2^-\right] \phi_3^+ \phi_4^-\notag\\[5pt]
        &\hspace{70pt}+\phi_1^- \phi_2^- \phi_3^- \phi_4^++ \left[\phi_3^-, \phi_4^+\right] \phi_1^- \phi_2^-+ \left[\phi_2^-, \phi_4^+\right] \phi_1^- \phi_3^-+ \left[\phi_2^-, \phi_3^- \right] \phi_1^- \phi_4^++\phi_1^- \phi_2^- \phi_3^- \phi_4^-\notag\\[5pt]
        &= :\phi_1 \phi_2 \phi_3 \phi_4: + \Delta_F(x_1 - x_2) \Delta_F(x_3 - x_4) + \Delta_F(x_1 - x_3) \Delta_F(x_2 - x_4) + \Delta_F(x_1 - x_4) \Delta_F(x_2 - x_3)
    \end{align*}
As a re-ordering of $x_1^0>x_2^0>x_3^0>x_4^0$ corresponds to a different choice of labels and since $\Delta_F$ and normal ordered products are symmetric, this calculation suffices to prove Wick’s theorem.

Thus, Wick's theorem is verified for the product of four scalar fields.